\documentclass[11pt]{article}
\usepackage[T1]{fontenc}
\usepackage{textcomp}
\usepackage[margin=0.75in]{geometry}
\usepackage{amsmath,amssymb,amsthm}
\usepackage{array,booktabs}
\usepackage{hyperref}
\setlength{\parindent}{0pt}
\newtheorem{theorem}{Theorem}
\theoremstyle{definition}
\newtheorem{definition}{Definition}
\title{Flow Proof Artifact: prop-12-gnomon-sum-squares}
\author{Generated by \texttt{flow doc proof}}
\date{Flow Proof Book}
\begin{document}
\maketitle
To apply a gnomon equal to the excess of one square over another in a parallelogram.\\[0.5em]
\textbf{Source.} euclid — Elements, Book II, Proposition 12\\[0.5em]
\bigskip
\noindent{\large \textbf{Derived fact 1.} \textit{To apply a gnomon equal to the excess of one square over another in a parallelogram} \hfill the Euclidean plane \cdot Euclid Book II \cdot Proposition 12: a gnomon equals the excess of two squares}\par
\smallskip\noindent\textit{Coordinate: the Euclidean plane \cdot Euclid Book II \cdot Proposition 12: a gnomon equals the excess of two squares}\par
\smallskip\noindent\textit{Source: euclid — Elements, Book II, Proposition 12}\par
\smallskip\noindent\textit{Built on: proposition 11: a gnomon equals the rectangle by the segments, for Euclid Book II on the Euclidean plane, proposition 4: the square on the whole equals the squares on the parts plus twice their rectangle, for Euclid Book II on the Euclidean plane}\par
\medskip
\noindent\textbf{Goal.} To apply a gnomon equal to the excess of one square over another in a parallelogram\par
\noindent$\displaystyle \text{gnomon equals sq A minus sq B}$\par
\medskip
\noindent
\begin{tabular}{@{} >{\raggedright\arraybackslash}p{0.06\textwidth} >{\raggedright\arraybackslash}p{0.47\textwidth} >{\raggedleft\arraybackslash}p{0.41\textwidth} @{}}
\textbf{\#} & \textbf{Proof} & \textbf{Mathematics} \\
\midrule
\textbf{1.} & We prove that proposition 12: a gnomon equals the excess of two squares for Euclid Book II the Euclidean plane. & \\[0.45em]
\textbf{2.} & Let sq\_A = square\_on\_larger. & \\[0.45em]
\textbf{3.} & Let sq\_B = square\_on\_smaller. & \\[0.45em]
\textbf{4.} & Let gnomon = gnomon\_in\_parallelogram. & \\[0.45em]
\textbf{5.} & From step 2, step 3, and step 4, we can deduce that gnomon equals sq A - sq B. & $\displaystyle gnomon = \text{sq A} - \text{sq B}$ \\[0.45em]
\textbf{6.} & We invoke the derived fact governing Euclid Book II the Euclidean plane: proposition 11: a gnomon equals the rectangle by the segments, for Euclid Book II on the Euclidean plane. & \\[0.45em]
\textbf{7.} & We invoke the derived fact governing Euclid Book II the Euclidean plane: proposition 4: the square on the whole equals the squares on the parts plus twice their rectangle, for Euclid Book II on the Euclidean plane. & \\[0.45em]
\textbf{8.} & From step 6 and step 7, this implies gnomon equals sq A minus sq B. Hence proven. & $\displaystyle \text{gnomon equals sq A minus sq B}$ \\[0.45em]
\bottomrule
\end{tabular}
\label{thm:Geometry-Euclid Book II-proposition_12_a_gnomon_equals_the_excess_of_two_squares}
\par\medskip\hrule
\end{document}
